% Frontera de truncamiento de la vía dodecafásica completa y delimitación
% exacta del resolvente analítico finito.
\paragraph{Dependencia exacta respecto de la frontera de truncamiento.}
\label{ampl-alfa-frontera-exacta}
Las coordenadas regionales y el registro dodecafásico proceden de la
generación común APP--TRIT--TPK. Una vez producidas, su composición
posicional posee una dependencia explícita respecto de la profundidad
de lectura. Sean $B=1000$,
\[
 Z_j=P_j+E_j-\Phi_j-K_j^{\mathrm{per}},\qquad
 y_N=\sum_{j=1}^{N}Z_jB^{-j},\qquad
 R_{N+1}(Z)=\sum_{r\geq0}Z_{N+1+r}B^{-r-1}.
\]
La serie completa de \eqref{eq:alpha-serie} satisface exactamente
\[
 \alpha_A=y_N+B^{-N}R_{N+1}(Z),\qquad
 -2<R_{N+1}(Z)<2.
\]
Esta identidad reúne la contribución de la ventana y la de su
prolongación, ambas expresadas en las mismas coordenadas.

\begin{proposition}[Transporte de la frontera y cilindro canónico]
\label{ampl-alfa-prop-frontera}
Sea $A_j^{(N,t)}$ la salida de la división posicional finita con
$c_{N+1}=t\in\mathcal C=\{-2,-1,0,1,2\}$, y sea
$a_N^{(t)}=\sum_{j=1}^{N}A_j^{(N,t)}B^{-j}$. Para los registros del
dominio considerado, se cumple
\begin{equation}
 \alpha_A-a_N^{(t)}
 =B^{-N}\bigl(R_{N+1}(Z)-t\bigr),
 \qquad
 |\alpha_A-a_N^{(t)}|<4B^{-N}.
 \label{ampl-alfa-eq-frontera}
\end{equation}
El cociente de la prolongación completa es
$c_{N+1}^{\infty}=\lfloor R_{N+1}(Z)\rfloor$. Con esta frontera,
\[
 0\leq\alpha_A-a_N^{(c_{N+1}^{\infty})}<B^{-N},
\]
de modo que la ventana obtenida es el prefijo canónico de orden $N$.
\end{proposition}
\begin{proof}
El telescopado de \eqref{eq:alpha-identidad-local} da
$a_N^{(t)}=y_N+tB^{-N}-c_1$. Como $Z_1=7$, cada estado entrante
pertenece a $\mathcal C$ y el primer dividendo queda entre $5$ y $9$,
se tiene $c_1=0$. Sustituir esta igualdad en la descomposición de la
serie prueba \eqref{ampl-alfa-eq-frontera}. La cota resulta de
$R_{N+1}(Z)\in(-2,2)$ y $t\in\mathcal C$.

Para la normalización canónica completa,
$c_{N+1}^{\infty}=R_{N+1}(Z)-R_{N+1}(A)$ y
$R_{N+1}(A)\in[0,1)$. La integralidad del cociente implica
$c_{N+1}^{\infty}=\lfloor R_{N+1}(Z)\rfloor$. La fórmula de error
se convierte entonces en
$\alpha_A-a_N^{(c_{N+1}^{\infty})}=B^{-N}R_{N+1}(A)$, que prueba
la pertenencia al cilindro canónico.
\end{proof}

La proposición identifica el significado de la frontera derecha: es la
parte entera de la cola firmada expresada en la escala de la posición
siguiente. El cálculo sobre las cinco fronteras conserva esta información
como una familia finita de posibilidades hasta que los bloques a la
izquierda se estabilizan. La contracción cuantitativa procede de la
profundidad y del transporte de la cola; su objeto es la normalización de
las coordenadas ya generadas.

\paragraph{Dominio graduado y estatuto de la representación analítica.}
\label{ampl-alfa-dominio-resolvente}
La composición analítica exhibida en
\S\ref{subsec:alpha-resolvente} actúa sobre
$\operatorname{span}\{e_7,e_8,e_9\}$ y satisface $Ne_9=0$. Su
evaluación produce exactamente $T_{7:9}$; el grado $9$ es la última
componente de ese dominio. La identidad $N^3=0$ explica por qué el
resolvente se expresa mediante tres sumandos. El resolvente local no determina
por sí mismo coeficientes de grado superior. La
sección~\ref{subsec:alpha-limite-dos-vias} demuestra esa no unicidad del jet
desnudo y separa dos operaciones: la publicación dodecafásica produce la
coordenada \(\alpha_{\mathrm{HMT}}\) desde los registros del estado; después,
la representación analítica consume la misma precoordenada y genera, por una
recurrencia explícita, todos los coeficientes de su imagen analítica. No se
consulta una raíz convencional, ni se promueve la representación analítica a
generador independiente. Los cuadrados de compatibilidad y los cilindros
anidados prueban la igualdad de ambas publicaciones a cada profundidad. La
nilpotencia local, la publicación dodecafásica y la representación analítica
desempeñan así funciones distintas: la primera certifica el resultado finito
de orden nueve; la segunda determina \(\alpha_{\mathrm{HMT}}\); la tercera
representa esa misma coordenada a toda profundidad y verifica su compatibilidad.
